\documentclass[11pt,a4paper]{article}
\usepackage[margin=1in]{geometry}
\usepackage{hyperref}
\usepackage{enumitem}

\title{Strategic Management \\ Week 10 \\ Slide-by-Slide Notes (ELI5) \\ Unit 4}
\author{(Generated from \texttt{sm-week-10-slides-5.pdf})}
\date{}

\begin{document}
\maketitle

\section*{How to use this (ELI5)}
\begin{itemize}
  \item For each slide, write the ``Exam write'' sentence in your notes.
  \item Then copy the ELI5 explanation to remember the logic.
  \item Focus on model meaning, not only keywords.
\end{itemize}

\section*{Slide 1 of 23 (Contents: Unit 4)}
\textbf{Key idea (from slide):} Internal analysis, value chain, resources and capabilities, SWOT.\\
\textbf{ELI5:} This unit teaches how to look inside the firm to find real strengths and weaknesses.\\
\textbf{Exam write:} ``Unit 4 is about internal diagnosis: value chain + RBV + SWOT.''

\section*{Slide 2 of 23 (Learning objectives)}
\textbf{Key idea (from slide):} Make strategic diagnosis, understand sustained competitiveness sources, analyze resources/capabilities, and connect external O/T with internal S/W for decisions.\\
\textbf{ELI5:} You learn what the company is really good at, and how to use that in strategy.\\
\textbf{Exam write:} ``Internal analysis supports strategic decisions by identifying sustainable strengths and key weaknesses.''

\section*{Slide 3 of 23 (Internal diagnosis tools)}
\textbf{Key idea (from slide):} Main tools: firm identity and strategic profile, value chain, resources and capabilities.\\
\textbf{ELI5:} First, understand who the company is; second, understand how it works; third, understand what assets it has.\\
\textbf{Exam write:} ``A complete internal diagnosis combines identity profile, value chain, and resources/capabilities analysis.''

\section*{Slide 4 of 23 (Firm identity: key variables)}
\textbf{Key idea (from slide):} Identity variables include age, scope, size, ownership, and activity type.\\
\textbf{ELI5:} A startup and a large multinational cannot use the same strategy style.\\
\textbf{Exam write:} ``Before strategy choice, classify firm identity (age, scope, size, ownership, activity).''

\section*{Slide 5 of 23 (Strategic profile by function)}
\textbf{Key idea (from slide):} Profile across operations, marketing, finance, people; compare to rival/industry average. Pros: simple/intuitive. Cons: subjective/lack of comparison precision.\\
\textbf{ELI5:} You score each area to see where you are strong or weak.\\
\textbf{Exam write:} ``Strategic profile is useful for quick diagnosis but must be validated because it can be subjective.''

\section*{Slide 6 of 23 (Value chain concept)}
\textbf{Key idea (from slide):}
\begin{itemize}
  \item Porter value chain = breakdown of activities needed to produce/deliver.
  \item Primary and supporting activities.
  \item Value system includes suppliers and customers.
  \item Goal: identify strengths/weaknesses, core competences, outsourcing options, and sources of SCA.
\end{itemize}
\textbf{ELI5:} Value chain shows where value is created and where money leaks.\\
\textbf{Exam write:} ``Use value chain to map activity-level strengths, weaknesses, and strategic advantage sources.''

\section*{Slide 7 of 23 (Porter value chain model)}
\textbf{Key idea (from slide):} Graphic of Porter (1985) value chain architecture.\\
\textbf{ELI5:} This is the map of how work flows inside the company.\\
\textbf{Exam write:} ``Reference Porter value chain structure when explaining cost or differentiation sources.''

\section*{Slide 8 of 23 (Sources of competitive advantage in value chain)}
\textbf{Key idea (from slide):}
\begin{itemize}
  \item Primary and support activities themselves
  \item Horizontal links inside the chain
  \item Vertical links with suppliers/buyers
  \item Through coordination and optimization
\end{itemize}
\textbf{ELI5:} Advantage comes not only from one activity, but from how activities connect.\\
\textbf{Exam write:} ``Competitive advantage can come from activities and from horizontal/vertical link coordination.''

\section*{Slide 9 of 23 (Exhibit 4.1 Ford vertical links)}
\textbf{Key idea (from slide):} Real example of vertical links in Ford Valencia factory.\\
\textbf{ELI5:} Supplier and buyer links can create efficiency and reliability.\\
\textbf{Exam write:} ``Use vertical-link examples to show how value system coordination supports advantage.''

\section*{Slide 10 of 23 (Illustration 4.1 Mercadona value chain/system)}
\textbf{Key idea (from slide):} Analyze Mercadona value creation in primary/support activities and identify horizontal/vertical coordination and optimization examples.\\
\textbf{ELI5:} Practice applying the model to a real company.\\
\textbf{Exam write:} ``In case answers, identify at least one horizontal and one vertical link with strategic effect.''

\section*{Slide 11 of 23 (RBV paradigm and aim)}
\textbf{Key idea (from slide):} Resource-Based View is current paradigm; aim is to discover and assess internal strengths (strategic resources/capabilities).\\
\textbf{ELI5:} RBV says firms differ because their internal assets and routines differ.\\
\textbf{Exam write:} ``RBV focuses on internal strategic resources and capabilities as roots of sustained advantage.''

\section*{Slide 12 of 23 (Resources, capabilities, environment, and SCA)}
\textbf{Key idea (from slide):} Resources (tangible/intangible/human) plus organizational capabilities lead to strategy and competitive advantage under environment opportunities/threats.\\
\textbf{ELI5:} The outside world gives pressure; inside assets decide if the firm can win.\\
\textbf{Exam write:} ``Competitive advantage emerges when resources and capabilities are matched to external opportunities/threats.''

\section*{Slide 13 of 23 (RBV key premises and steps)}
\textbf{Key idea (from slide):}
\begin{itemize}
  \item Key premises: firm heterogeneity and imperfect mobility of resources/capabilities.
  \item Steps: identify R\&C, evaluate role for sustained advantage, formulate corporate/business strategy for exploitation/improvement.
\end{itemize}
\textbf{ELI5:} Firms are not equal, and many resources cannot be moved easily, so internal differences matter.\\
\textbf{Exam write:} ``RBV process: identify resources -> evaluate strategic value -> formulate strategy for exploitation and development.''

\section*{Slide 14 of 23 (Resources concept and categories)}
\textbf{Key idea (from slide):} Resources are assets the firm controls; types include tangible (financial/material) and intangible (human-based, non-human, organizational, technological).\\
\textbf{ELI5:} Resources are the ingredients the firm owns or controls.\\
\textbf{Exam write:} ``Classify resources clearly before claiming strategic importance.''

\section*{Slide 15 of 23 (Capabilities concept)}
\textbf{Key idea (from slide):} Capabilities are coordination/combination of resources; include core/distinctive competences and organizational routines; types include cultural, functional, dynamic.\\
\textbf{ELI5:} Capability is how well the firm combines resources to do things better than rivals.\\
\textbf{Exam write:} ``Capabilities are coordinated routines that transform resources into competitive outcomes.''

\section*{Slide 16 of 23 (Illustration 4.2 Intangibles and capabilities as SCA sources)}
\textbf{Key idea (from slide):} Critical thinking prompt on why intangible resources and capabilities are main SCA sources.\\
\textbf{ELI5:} Intangibles (brand, know-how, culture) are usually harder to copy than machines.\\
\textbf{Exam write:} ``Intangible resources and capabilities often drive SCA because imitation is difficult and slow.''

\section*{Slide 17 of 23 (Strategic analysis of R\&C: three goals)}
\textbf{Key idea (from slide):} Not all resources/capabilities create SCA. Strategic analysis tests: generating advantage, sustaining advantage, and rent appropriation.\\
\textbf{ELI5:} A resource is strategic only if it helps win now, keep winning later, and keep the value inside the firm.\\
\textbf{Exam write:} ``Evaluate R\&C by generation, sustainability, and appropriation of value/rents.''

\section*{Slide 18 of 23 (Generation criteria: scarcity and relevance)}
\textbf{Key idea (from slide):}
\begin{itemize}
  \item Scarcity, not available to many competitors.
  \item Relevance, overlap with strategic industry factors.
\end{itemize}
\textbf{ELI5:} If everyone has it, it is not an advantage; if it does not matter in this industry, it is not strategic.\\
\textbf{Exam write:} ``For advantage generation, test scarcity and relevance to key industry success factors.''

\section*{Slide 19 of 23 (Sustainability criteria)}
\textbf{Key idea (from slide):}
\begin{itemize}
  \item Durability
  \item Inimitability
  \item Low tradeability
  \item Limited substitutability
  \item Complementarity
\end{itemize}
\textbf{ELI5:} Advantage lasts when rivals cannot copy, buy, or replace your key resources easily.\\
\textbf{Exam write:} ``SCA sustainability needs durability + imitation barriers + low tradeability/substitutability + complementarity.''

\section*{Slide 20 of 23 (Rent appropriation)}
\textbf{Key idea (from slide):} Appropriability = control rights over strategic R\&C and resulting rents; appropriation differs by resource type (intangible often harder to define/control, tangible easier).\\
\textbf{ELI5:} Winning is not enough; you must also keep the value your strategy creates.\\
\textbf{Exam write:} ``Appropriability asks who captures value from strategic resources and under what control rights.''

\section*{Slide 21 of 23 (Illustration 4.3 Inditex strategic analysis)}
\textbf{Key idea (from slide):} Exercise to identify/measure two resources (tangible + intangible) and one key capability in Inditex.\\
\textbf{ELI5:} Practice turning theory into concrete company evidence.\\
\textbf{Exam write:} ``In firm cases, identify specific resources and one capability, then evaluate strategic quality.''

\section*{Slide 22 of 23 (Managing resources and capabilities)}
\textbf{Key idea (from slide):}
\begin{itemize}
  \item Exploit R\&C internally and externally.
  \item Improve R\&C provision via external acquisition and internal development.
\end{itemize}
\textbf{ELI5:} Firms must use current strengths and also build new ones.\\
\textbf{Exam write:} ``R\&C strategy includes both exploitation of existing assets and development/acquisition of new assets.''

\section*{Slide 23 of 23 (SWOT summary and limits)}
\textbf{Key idea (from slide):}
\begin{itemize}
  \item SWOT summarizes Opportunities/Threats and Strengths/Weaknesses.
  \item Pros: popular, overall view, simple.
  \item Cons: static, weak integration, no direct strategy recipe.
  \item Use logic: maintain strengths, correct weaknesses, avoid threats, exploit opportunities.
\end{itemize}
\textbf{ELI5:} SWOT is a good summary page, but not a full strategy by itself.\\
\textbf{Exam write:} ``Use SWOT as diagnosis synthesis, then add strategic choices and implementation logic.''

\end{document}

